\chapter*{Preface}
\addcontentsline{toc}{chapter}{Preface}
% STATUS: written 2026-11-06 (states the scope, the source rule and the living
% edition rule; these are the book's own policies, fixed in CLAUDE.md and in the
% ROADMAP, not external claims.)
\label{page:preface}

A derivation is the algebraic form of a derivative.
Leibniz wrote the product rule in 1684, and three centuries later the same rule
defines a map on any ring: an additive map $d$ with $d(xy)=d(x)y+xd(y)$.
Every chapter of this book studies some version of that equation, either the
equation itself or a weaker or stronger form of it.

\section*{Why a encyclopaedia and not a survey}

A survey selects and argues. A encyclopaedia tries to be complete in scope, so that a
reader can find where a given question is treated and what is known about it.

This book is exhaustive in scope, not in text.
Every flavour of derivation, every class of ring and algebra where derivations
act, and every line of development from Leibniz to this year has a place in the
plan.
The plan is fixed in advance and the text grows into it.

\section*{The source rule}

No fact, definition, example or application is stated before its source has been
read in full from the page.
A recalled theorem appears as a theorem-like environment with its located
citation, keeping the kind and the wording of the source.
Proofs and remarks of the book itself are marked as its own.

This rule has a cost, and the cost is visible.
A section whose sources are not yet obtained stays a marked planned
stub, and the stub says which source it waits for.
A reader who finds a stub knows that the coverage is planned there.
The text is not yet written, which is a different thing from a section that was
written and then cut.

This edition has no such stub.
Every section is written.
Some sources cited in the bibliography could not be opened at all.
Their entries carry the marker `unverified-ref` together with the reason, and
their sections cite them for what the record states rather than quoting a
theorem from a page nobody has seen.
The count is kept in the download record of this book, not here, because it
changes as copies are found.
The work list that the search step keeps records the same fact for the works
that were found and not cited: a lead is not a source until its page has been
read.
The policy stays in place for later editions, which is why it is stated here.

\section*{Where the applications are}

A derivation is a rule that obeys the Leibniz product rule. The rule itself
goes back to Leibniz \cite{Leibniz1684}. The name was already familiar to
Ritt when he built differential algebra on it \cite[p.~1]{Ritt1950}. The same
rule appears again in nine fields beyond the differential calculus. The nine
names are classical and quantum mechanics, differential geometry,
differential algebra, coding, control, group analysis, model theory,
algebraic topology and automatic differentiation. Each field is reported in
the chapter that owns its notions.
Each report names one field, cites the source that names the derivation it
rests on, and reads the definition or theorem of that source. The reports do
not narrate: they record.

Physics was the first field that asked why a derivation arises. Dirac
\cite{Dirac1925} translated noncommutative kinematics into the language of a
Poisson bracket
(Section~\ref{sec:basic-theory:physics}), and Kadison \cite{Kadison1966} later
asked whether every derivation of an operator algebra comes from inside
(Section~\ref{sec:banach-operator:kadison-sakai}). Geometry came next: a tangent vector
at a point of a smooth manifold is a derivation of the smooth functions there
(Section~\ref{sec:commutative:geometry}). Differential algebra began when Ritt
and Kolchin axiomatized the symbol $d/dx$ as a derivation on a field, and the
field-of-definition question for the Painlev\'e equations sits inside that
axiomatization (Section~\ref{sec:differential-algebra:differential-equations}). In
positive characteristic the ordinary product rule survives in the
divided-power form of Hasse and Schmidt; the higher derivations are the
algebraic backbone of Reed--Solomon and Goppa codes
(Section~\ref{sec:higher-multi:coding}). Control theory discovered that an
input--output system is a finitely generated differential field extension,
with the input a differential transcendence basis
(Section~\ref{sec:differential-algebra:control}). On a group the same rule is a
derivation of the group ring, and Fox showed the free group carries one
derivation per generator (Section~\ref{sec:group-hopf:harmonic-analysis}).
Ax translated the conjecture of Schanuel about the complex exponential into a
statement about a derivation on a differential field. His 1971 theorem is the
seed of the Ax--Schanuel inequalities
(Section~\ref{sec:differential-algebra:transcendence}). Hochschild and cyclic
homology compute derivations as the homology of a chain complex, and on a
smooth commutative algebra the answer is the exterior algebra of K\"ahler
differentials (Section~\ref{sec:cohomology:algebraic-topology}). A computer
program that differentiates a function runs the chain rule and the Leibniz
rule step by step under automatic differentiation
(Section~\ref{sec:differential-algebra:automatic-differentiation}).

\hbfigure{02-basic-theory/oscillator}{The harmonic oscillator. The derivation $f\mapsto\{f,H\}$ with $H=(p^2+q^2)/2$ is the vector field $(p,-q)$ on the phase plane; its flow is a rotation.}{fig:basic:oscillator}

\section*{The living edition}

The subject grows.
Papers on derivations appear every month, and a printed encyclopaedia would be out of
date before it reached a reader.
This edition is therefore updated monthly: new work is added where it changes a
statement, a table entry or an open problem.
Appendix~\ref{app:atlas} reports what the edition has read, and the version
number on the title page names the edition.

\section*{How the book was made}

The work of reading, checking and drafting was carried out with artificial
intelligence under human direction.
A person set the scope, fixed the questions and decided what the book would
contain.
Machine tools assisted the reading of each source, the extraction of its results
and the drafting of each section.
Every claim was then checked against the page it cites.
The rules that govern this are stated in the repository of the book.
The records of what was read are kept there too, so a reader can trace any
statement to its source.

The name on the title page is the person who directed the work and answers for
it.
